% GENERATED FILE. Edit canonical lessons, then run npm run generate:books.
% canonical-source-hash: fnv1a64:8f534454b84ba05a
% canonical-lessons: SA-C430-hattah, SA-C430-vikreta, SA-C430-kreta, SA-C430-grahakah, SA-C430-panate

\chapter{At the Bazaar}
\label{ch:sa-at-the-bazaar}
\hlchaptermodality{430}

\begin{chapteropening}
\textbf{By the end of this chapter:} \emph{I can talk with sellers and buyers at a bazaar.}

Complete the last lesson of chapter 430: I can talk with sellers and buyers at a bazaar.
\end{chapteropening}

\section[\texorpdfstring{\sk{हट्टः} (haṭṭaḥ)}{हट्टः (haṭṭaḥ)}]{\sk{हट्टः} (haṭṭaḥ) — a bazaar; a market fair}

\label{lesson:SA-C430-hattah}



\begin{quote}
Before the new one: say the Sanskrit for a door bolt, then the Sanskrit for rent.
\end{quote}

\subsection*{You'll want to know: \sk{हट्टः}}
\textbf{\sk{हट्टः}} — \emph{haṭṭaḥ} — \textquotedblleft{}a bazaar; a market fair\textquotedblright{}.

\begin{quote}
\sk{सोमवासरे} \sk{ग्रामे} \sk{हट्टः} \sk{भवति।}
\end{quote}

\begin{quote}
\emph{somavāsare grāme haṭṭaḥ bhavati.}
\end{quote}

\begin{quote}
On Monday there is a bazaar in the village.
\end{quote}

A masculine noun.

\begin{cousinweb}
\textbf{\sk{हट्टः}} (\emph{haṭṭaḥ}) is the word MW gives for a market, a fair; MW gives no derivation, so its origin is uncertain; it lives on as Hindi \emph{hāṭ}, \textquotedblleft{}a market\textquotedblright{} (MW; CDIAL 13944).
\end{cousinweb}

\begin{grammarlens}[title={What you've built}]
One more word for getting around.
\end{grammarlens}

\subsection*{Your turn}
\begin{itemize}
\raggedright
  \item \emph{Say it:} \emph{haṭṭaḥ}
  \item \emph{Say it:} \emph{haṭṭaḥ}, once more
  \item \emph{Say it:} \emph{somavāsare grāme haṭṭaḥ bhavati.} — On Monday there is a bazaar in the village.
  \item \emph{Recall:} say the Sanskrit for a door bolt, then the Sanskrit for rent, then say \emph{haṭṭaḥ} again
\end{itemize}

\begin{tcolorbox}[breakable,colback=teal!4,colframe=teal!35!black,title={Before you move on}]
What does \sk{हट्टः} mean? (\textquotedblleft{}A bazaar; a market fair\textquotedblright{}.) Say it once more.
\end{tcolorbox}
\section[\texorpdfstring{\sk{विक्रेता} (vikretā)}{विक्रेता (vikretā)}]{\sk{विक्रेता} (vikretā) — a seller}

\label{lesson:SA-C430-vikreta}



\begin{quote}
Before the new one: say the Sanskrit for rent, then the Sanskrit for a bazaar.
\end{quote}

\subsection*{You'll want to know: \sk{विक्रेता}}
\textbf{\sk{विक्रेता}} — \emph{vikretā} — \textquotedblleft{}a seller\textquotedblright{}.

\begin{quote}
\sk{विक्रेता} \sk{उच्चैः} \sk{वदति।}
\end{quote}

\begin{quote}
\emph{vikretā uccaiḥ vadati.}
\end{quote}

\begin{quote}
The seller calls out loudly.
\end{quote}

A masculine noun.

\begin{cousinweb}
\textbf{\sk{विक्रेता}} (\emph{vikretā}) is built from \emph{vi} and the root √krī, \textquotedblleft{}to buy\textquotedblright{}, which with \emph{vi} means to sell, as in \emph{vikrīṇāti}, \textquotedblleft{}sells\textquotedblright{}, with the agent ending \emph{-tṛ}; MW gives \emph{vikretṛ} as a seller (MW).
\end{cousinweb}

\begin{grammarlens}[title={What you've built}]
One more word for getting around.
\end{grammarlens}

\subsection*{Your turn}
\begin{itemize}
\raggedright
  \item \emph{Say it:} \emph{vikretā}
  \item \emph{Say it:} \emph{vikretā}, once more
  \item \emph{Say it:} \emph{vikretā uccaiḥ vadati.} — The seller calls out loudly.
  \item \emph{Recall:} say the Sanskrit for rent, then the Sanskrit for a bazaar, then say \emph{vikretā} again
\end{itemize}

\begin{tcolorbox}[breakable,colback=teal!4,colframe=teal!35!black,title={Before you move on}]
What does \sk{विक्रेता} mean? (\textquotedblleft{}A seller\textquotedblright{}.) Say it once more.
\end{tcolorbox}
\section[\texorpdfstring{\sk{क्रेता} (kretā)}{क्रेता (kretā)}]{\sk{क्रेता} (kretā) — a buyer}

\label{lesson:SA-C430-kreta}



\begin{quote}
Before the new one: say the Sanskrit for a bazaar, then the Sanskrit for a seller.
\end{quote}

\subsection*{You'll want to know: \sk{क्रेता}}
\textbf{\sk{क्रेता}} — \emph{kretā} — \textquotedblleft{}a buyer\textquotedblright{}.

\begin{quote}
\sk{क्रेता} \sk{हट्टे} \sk{मूल्यं} \sk{पृच्छति।}
\end{quote}

\begin{quote}
\emph{kretā haṭṭe mūlyaṃ pṛcchati.}
\end{quote}

\begin{quote}
The buyer asks the price at the bazaar.
\end{quote}

A masculine noun.

\begin{cousinweb}
\textbf{\sk{क्रेता}} (\emph{kretā}) is built on the root √krī, \textquotedblleft{}to buy\textquotedblright{}, which also gives \emph{krīṇāti}, \textquotedblleft{}buys\textquotedblright{}, with the agent ending \emph{-tṛ}; MW gives \emph{kretṛ} as a buyer, a purchaser (MW).
\end{cousinweb}

\begin{grammarlens}[title={What you've built}]
One more word for getting around.
\end{grammarlens}

\subsection*{Your turn}
\begin{itemize}
\raggedright
  \item \emph{Say it:} \emph{kretā}
  \item \emph{Say it:} \emph{kretā}, once more
  \item \emph{Say it:} \emph{kretā haṭṭe mūlyaṃ pṛcchati.} — The buyer asks the price at the bazaar.
  \item \emph{Recall:} say the Sanskrit for a bazaar, then the Sanskrit for a seller, then say \emph{kretā} again
\end{itemize}

\begin{tcolorbox}[breakable,colback=teal!4,colframe=teal!35!black,title={Before you move on}]
What does \sk{क्रेता} mean? (\textquotedblleft{}A buyer\textquotedblright{}.) Say it once more.
\end{tcolorbox}
\section[\texorpdfstring{\sk{ग्राहकः} (grāhakaḥ)}{ग्राहकः (grāhakaḥ)}]{\sk{ग्राहकः} (grāhakaḥ) — a customer}

\label{lesson:SA-C430-grahakah}



\begin{quote}
Before the new one: say the Sanskrit for a seller, then the Sanskrit for a buyer.
\end{quote}

\subsection*{You'll want to know: \sk{ग्राहकः}}
\textbf{\sk{ग्राहकः}} — \emph{grāhakaḥ} — \textquotedblleft{}a customer\textquotedblright{}.

\begin{quote}
\sk{सायम्} \sk{आपणे} \sk{ग्राहकाः} \sk{आगच्छन्ति।}
\end{quote}

\begin{quote}
\emph{sāyam āpaṇe grāhakāḥ āgacchanti.}
\end{quote}

\begin{quote}
In the evening customers come to the shop.
\end{quote}

A masculine noun.

\begin{cousinweb}
\textbf{\sk{ग्राहकः}} (\emph{grāhakaḥ}) is built on the root √grah, \textquotedblleft{}to take\textquotedblright{}, which also gives \emph{gṛhṇāti}, \textquotedblleft{}takes\textquotedblright{}, with the agent ending \emph{-aka}; MW gives \emph{grāhaka} as one who takes and as a purchaser (MW).
\end{cousinweb}

\begin{grammarlens}[title={What you've built}]
One more word for getting around.
\end{grammarlens}

\subsection*{Your turn}
\begin{itemize}
\raggedright
  \item \emph{Say it:} \emph{grāhakaḥ}
  \item \emph{Say it:} \emph{grāhakaḥ}, once more
  \item \emph{Say it:} \emph{sāyam āpaṇe grāhakāḥ āgacchanti.} — In the evening customers come to the shop.
  \item \emph{Recall:} say the Sanskrit for a seller, then the Sanskrit for a buyer, then say \emph{grāhakaḥ} again
\end{itemize}

\begin{tcolorbox}[breakable,colback=teal!4,colframe=teal!35!black,title={Before you move on}]
What does \sk{ग्राहकः} mean? (\textquotedblleft{}A customer\textquotedblright{}.) Say it once more.
\end{tcolorbox}
\section[\texorpdfstring{\sk{पणते} (paṇate)}{पणते (paṇate)}]{\sk{पणते} (paṇate) — haggles; bargains}

\label{lesson:SA-C430-panate}



\begin{quote}
Before the new one: say the Sanskrit for a buyer, then the Sanskrit for a customer.
\end{quote}

\subsection*{You'll want to know: \sk{पणते}}
\textbf{\sk{पणते}} — \emph{paṇate} — \textquotedblleft{}haggles; bargains\textquotedblright{}.

\begin{quote}
\sk{ग्राहकः} \sk{हट्टे} \sk{पणते।}
\end{quote}

\begin{quote}
\emph{grāhakaḥ haṭṭe paṇate.}
\end{quote}

\begin{quote}
The customer haggles at the bazaar.
\end{quote}

Past (he …): \sk{सः} \sk{हट्टे} \sk{पणितवान्}, \emph{saḥ haṭṭe paṇitavān}.

\begin{cousinweb}
\textbf{\sk{पणते}} (\emph{paṇate}) comes from the root √paṇ, \textquotedblleft{}to barter, to bargain\textquotedblright{}, which also gives \emph{paṇa}, \textquotedblleft{}a stake, a coin\textquotedblright{}; MW gives the participle \emph{paṇita}, \textquotedblleft{}transacted as business\textquotedblright{} (MW).
\end{cousinweb}

\begin{grammarlens}[title={What you've built}]
One more word for getting around.
\end{grammarlens}

\subsection*{Your turn}
\begin{itemize}
\raggedright
  \item \emph{Say it:} \emph{paṇate}
  \item \emph{Say it:} \emph{paṇate}, once more
  \item \emph{Say it:} \emph{grāhakaḥ haṭṭe paṇate.} — The customer haggles at the bazaar.
  \item \emph{Recall:} say the Sanskrit for a buyer, then the Sanskrit for a customer, then say \emph{paṇate} again
\end{itemize}

\begin{tcolorbox}[breakable,colback=teal!4,colframe=teal!35!black,title={Before you move on}]
What does \sk{पणते} mean? (\textquotedblleft{}Haggles; bargains\textquotedblright{}.) Say it once more.
\end{tcolorbox}
